\section{Introduction}

The Russian mathematician Peterson was a student of Minding's, who
in turn was interested in deformations (through bending) of
surfaces\footnote{See Peterson's biography at \url{http://www-history.mcs.st-and.ac.uk/Biographies/Peterson.html}.}, but unfortunately most of his works
(including his independent discovery of the Codazzi--Mainardi
equations and of the Gau\ss--Bonnet theorem) were made known to
Western Europe mainly after they were translated in 1905 from
Russian to French (as is the case with his deformations of
quadrics~\cite{P}, originally published in 1883 in Russian).
Peterson's work on deformations of general quadrics preceded that
of Bianchi, Calapso, Darboux, Guichard and \c{T}i\c{t}eica's from
the years 1899--1906 by two decades; in particu\-lar Peterson's
$1$-dimensional family of deformations of surfaces admitting a~common {\it conjugate system}~$(u,v)$ (that is the second
fundamental form is missing mixed terms $du\odot dv$) are {\it
associates} (a~notion naturally appearing in the inf\/initesimal
deformation problem) to Bianchi's $1$-dimensional family of
surfaces satisfying $(\log K)_{uv}=0$ in the common asymptotic
coordinates~$(u,v)$, $K$~being the Gau\ss\ curvature (see Bianchi~\cite[Vol.~2, \S\S~294, 295]{B1}).

The work of these illustrious geometers on deformations in
$\mathbb{C}^3$ of quadrics in $\mathbb{C}^3$ (there is no other
class of surfaces for which an interesting theory of deformation
has been built) is one of the crowning achievements of the golden
age of classical geometry of surfaces and at the same time it
opened new areas of research (such as af\/f\/ine and projective
dif\/ferential geometry) continued later by other illustrious
geometers (Blaschke, Cartan, etc.).

Peterson's $1$-dimensional family of deformations of
$2$-dimensional quadrics is obtained by imposing an ansatz
naturally appearing from a geometric point of view, namely the
constraint that the common conjugate system of curves is given by
intersection with planes through an axis and tangent cones
centered on that axis; thus this result of Koenigs (see Darboux
\cite[\S\S~91--101]{D}) was again (at least when the cones are
tangent along plane curves) previously known to Peterson. Note
also that Calapso in \cite{Ca} has put Bianchi's B\"{a}cklund
transformation of deformations in~$\mathbb{C}^3$ of general
$2$-dimensional quadrics with center in terms of common conjugate
systems (the condition that the conjugate system on a
$2$-dimensional quadric is a conjugate system on one of its
deformations in~$\mathbb{C}^3$ was known to Calapso for a decade,
but the B\"{a}cklund transformation for general quadrics eluded
Calapso since the common conjugate system was best suited for this
transformation only at the analytic level).

Although this is the original approach Peterson used to f\/ind his
deformations of quadrics, other features of his approach will make
it amenable to higher dimensional generalizations, namely the
warping of linear element (the warping of the linear element of a
plane curve to get the linear element of a surface of revolution
$(d(f\cos(u^1)))^2+(d(f\sin(u^1)))^2=(df)^2+f^2(du^1)^2$ for
$f=f(u^2)$ is such an example) and separation of variables;
post-priori the common conjugate system property may be given a
geometric explanation analogous to that in dimension $3$.

In 1919--1920 Cartan has shown in~\cite{C} (using mostly projective
arguments and his exterior dif\/ferential systems in involution and
exteriorly orthogonal forms tools) that space forms of dimension
$n$ admit rich families of deformations (depending on $n(n-1)$
functions of one variable) in surrounding $(2n-1)$-dimensional
space forms, that such deformations admit lines of curvature
(given by a canonical form of exteriorly orthogonal forms; thus
they have f\/lat normal bundle; since the lines of curvature on
$n$-dimensional space forms (when they are considered by
def\/inition as quadrics in surrounding $(n+1)$-dimensional space
forms) are undetermined, the lines of curvature on the deformation
and their corresponding curves on the quadric provide the common
conjugate system) and that the co-dimension $(n-1)$ cannot be
lowered without obtaining rigidity as the deformation being the
def\/ining quadric.

In 1983 Berger, Bryant and Grif\/f\/iths~\cite{B} proved (including by
use of tools from algebraic geometry) in particular that Cartan's
essentially projective arguments (including the exterior part of
his exteriorly orthogonal forms tool) can be used to generalize
his results to $n$-dimensional general quadrics with positive
def\/inite linear element (thus they can appear as quadrics in
$\mathbb{R}^{n+1}$ or as space-like quadrics in
$\mathbb{R}^n\times(i\mathbb{R})$) admitting rich families of
deformations (depending on $n(n-1)$ functions of one variable) in
surrounding Euclidean space $\mathbb{R}^{2n-1}$, that the
co-dimension $(n-1)$ cannot be lowered without obtaining rigidity
as the deformation being the def\/ining quadric and that quadrics
are the only Riemannian $n$-dimensional manifolds that admit a
family of deformations in $\mathbb{R}^{2n-1}$ as rich as possible
for which the exteriorly orthogonal forms tool (naturally
appearing from the Gau\ss\ equations) can be applied.

Although Berger, Bryant and Grif\/f\/iths~\cite{B} do not explicitly
state the common conjugate system property (which together with
the non-degenerate joined second fundamental forms assumption
provides a tool similar to the canonical form of exteriorly
orthogonal forms), this will turn out to be the correct tool of
dif\/ferential geometry needed to attack the deformation problem for
higher dimensional quadrics; also at least for diagonal quadrics
without center Peterson's deformations of higher dimensional
quadrics will turn out to be amenable to explicit computations of
their B\"{a}cklund transformation\footnote{See Dinc\u{a} I.I., Bianchi's
B\"{a}cklund transformation for higher dimensional quadrics,
\href{http://arxiv.org/abs/0808.2007}{arXiv:0808.2007}.}.

All computations are local and assumed to be valid on their open
domain of validity without further details; all functions have the
assumed order of dif\/ferentiability and are assumed to be
invertible, non-zero, etc when required (for all practical
purposes we can assume all functions to be analytic).

Here we have the two main theorems concerning the
$(n-1)$-dimensional family of deformations of higher dimensional
general quadrics and respectively its maximality:

\begin{theorem}\label{theorem1}
The quadric
\[
\sum_{j=0}^n\frac{(x_j)^2}{a_j}=1,\qquad
a_j\in\mathbb{C}^*
\] distinct parameterized with the conjugate
system $(u^1,\dots ,u^n)\subset\mathbb{C}^n$ given by the spherical
coordinates on the unit sphere
$\mathbb{S}^n\subset\mathbb{C}^{n+1}$:
\begin{gather*}
\mathcal{X}=\sqrt{a_0}\mathbf{C}_0e_0+
\sum_{k=1}^n\sqrt{a_k}\mathbf{C}_k\sin\big(u^k\big)e_k,\qquad
\mathbf{C}_k:=\prod_{j=k+1}^n\cos(u^j)
\end{gather*}
and the sub-manifold
\begin{gather*}
\mathcal{X}_{\mathbf{z}}=\sum_{k=1}^{n-1}\mathbf{C}_kf_k\big(\mathbf{z},u^k\big)
\big(\cos\big(g_k\big(\mathbf{z},u^k\big)\big)e_{2k-2}
+\sin\big(g_k\big(\mathbf{z},u^k\big)\big)e_{2k-1}\big)+h\big(\mathbf{z},u^n\big)e_{2n-2}
\end{gather*}
 of
$\mathbb{C}^{2n-1}$ depending on the parameters
$\mathbf{z}=(z_1,z_2,\dots ,z_{n-1})\in\mathbb{C}^{n-1}$, $z_0:=1$
and with
\begin{gather}
f_k\big(z_{k-1},z_k,u^k\big):=\sqrt{(z_{k-1}-z_k)a_0+(a_k-z_{k-1}a_0)\sin^2(u^k)},\qquad k=1,\dots ,n-1,\nonumber\\
g_k\big(z_{k-1},z_k,u^k\big):=\int_0^{u^k}\frac{\sqrt{(z_{k-1}-z_k)a_0a_k+(a_k-z_{k-1}a_0)z_ka_0\sin^2(t)}}
{(z_{k-1}-z_k)a_0+(a_k-z_{k-1}a_0)\sin^2(t)}dt,\nonumber\\
h\big(z_{n-1},u^n\big):=\int_0^{u^n}\sqrt{a_n-(a_n-z_{n-1}a_0)\sin^2(t)}dt\label{eq:fgn}
\end{gather}
have the same linear element
$|d\mathcal{X}|^2{=}|d\mathcal{X}_{\mathbf{z}}|^2$. For
$z_1={\cdots}=z_{n-1}=0$ we get \mbox{$g_2={\cdots}=g_{n-1}=0$},
$\mathcal{X}=\mathcal{X}_{\mathbf{0}}$ with
$\mathbb{C}^{n+1}\hookrightarrow\mathbb{C}^{2n-1}$ as
$(x_0,x_1,\dots ,x_n)\mapsto(x_0,x_1,x_2,0,x_3,0,\dots ,x_{n-1},0,x_n)$.
For $z_1=\dots =z_{n-1}=1$ we get
$\mathcal{X}_{\mathbf{1}}=(\mathbf{x}_0,\dots ,\mathbf{x}_{2n-2})$
given by Peterson's formulae
\begin{gather}
\sqrt{(\mathbf{x}_{2k-2})^2+(\mathbf{x}_{2k-1})^2}=\sqrt{a_k-a_0}\mathbf{C}_k\sin\big(u^k\big),\nonumber\\
 \tan^{-1}\left(\frac{\mathbf{x}_{2k-1}}{\mathbf{x}_{2k-2}}\right)
=\frac{\sqrt{a_0}}{\sqrt{a_k-a_0}}\tanh^{-1}\big(\cos\big(u^k\big)\big), \qquad k=1,\dots,n-1,
\nonumber\\
\mathbf{x}_{2n-2}=\int_0^{u^n}\sqrt{a_n-(a_n-a_0)\sin^2(t)}dt.\label{eq:pfn}
\end{gather}
Moreover $(u^1,\dots ,u^n)$ form a conjugate system on
$\mathcal{X}_{\mathbf{z}}$ with non-degenerate joined second
fundamental forms $($that is $[d^2\mathcal{X}^TN\ \
d^2\mathcal{X}_{\mathbf{z}}^TN_{\mathbf{z}}]$ is a symmetric
quadratic $\mathbb{C}^n$-valued form which contains only
$(du^j)^2$ terms for $N$ normal field of $\mathcal{X}$ and
$N_{\mathbf{z}}=[N_1\   \dots   \ N_{n-1}]$ normal frame of~$\mathcal{X}_{\mathbf{z}}$ and the dimension $n$ cannot be lowered
for $\mathbf{z}$ in an open dense set$)$.
\end{theorem}
